\documentclass[10pt,a4paper]{amsart}
\usepackage[margin=19mm]{geometry}
\usepackage{amsmath,amssymb,mathtools}
\usepackage[scr=rsfs,cal=euler]{mathalfa}
\usepackage{xfrac}
\usepackage[unicode,hidelinks,bookmarks=false]{hyperref}
\newcommand{\ie}{i.e.\ }
\newcommand{\cf}{cf.\ }
\newcommand{\resp}{respectively\ }
\newcommand{\Subsection}{\subsection}
\providecommand{\nobreakdash}{\nobreak\hbox{-}}
\setlength{\emergencystretch}{2em}
\multlinegap0pt
\usepackage{bm}
\usepackage{bbm}
\usepackage{dsfont}
\usepackage{xspace}
\usepackage{cancel}
\usepackage{mathtools}
\usepackage{amsthm}
\makeatletter
\@ifundefined{LEM}{\newtheorem{LEM}{Lemma}}{}
\@ifundefined{PROP}{\newtheorem{PROP}[LEM]{Proposition}}{}
\makeatother
\newcommand{\Hom}{\mathop{\textrm{Hom}}\nolimits}
\newcommand{\bcdot}{{\mbox{\boldmath{$\cdot$}}}}
\newcommand{\dD}{{\mathcal{D}}}
\newcommand{\f}{\varphi}
\newcommand{\gge}{\mbox{\raisebox{0.25ex}{${\scriptscriptstyle\ge}$}}}
\newcommand{\kk}{\mathbbm{k}}
\newcommand{\lle}{\mbox{\raisebox{0.25ex}{${\scriptscriptstyle\le}$}}}
\newcommand{\tr}{{\mbox{$t$-struc}\-tu\-r}}
\title[Highest weight categories via pairs of dual exceptional sequ]{Highest weight categories via pairs of dual exceptional sequences}
\author{Agnieszka Bodzenta, Alexey Bondal}
\date{Source: arXiv:2601.22004v1 (CC BY 4.0)}
\begin{document}
\maketitle

\noindent\emph{Source and licence.} Highest weight categories via pairs of dual exceptional sequences. Version arXiv:2601.22004v1 (CC BY 4.0), distributed under the Creative Commons Attribution 4.0 International licence, as recorded in the arXiv licence metadata for this exact version and shown on the abstract page https://arxiv.org/abs/2601.22004v1. Full rights notice: licenses/arxiv-2601.22004-CC-BY-4.0.txt.\par\smallskip

\noindent\textit{Editorial scope.} Counted unit: the Prop whose source label is \texttt{None} in the article (source0.tex lines 752--758, source label \texttt{None}), delivered together with the article's own complete original proof of it (source0.tex lines 759--798, one \texttt{proof} environment of 40 lines). No mathematics is written, repaired or re-derived: statement and proof are the article's own text line for line. Required context retained uncounted: lem\_glued\_t-str (lines 416--426); eqtn\_relations (lines 725--727); eqtn\_L\_S\_3\_P\_2 (lines 736--739). Every definition, display and label of those lines is retained unchanged. Omitted from this package and not redistributed: 1--415, 427--724, 728--735, 740--751, 799--1478. Nothing inside a retained excerpt is omitted or altered: every retained body is delivered exactly as the article's lines are written, so no omission receipt of any kind accompanies this package. Citations: the retained excerpts carry no citation command at all, so no bibliography entry of the article is transcribed and no reference is redistributed. Reference adaptation: none was needed and none was made; every reference inside the delivered unit resolves inside it, so no unresolved reference or citation is produced. Printed slips kept literal: no printed word, symbol, hypothesis or number of the counted statement or of its proof was altered. Layout and class: the article's own class is \texttt{amsart} with packages including amsfonts, amsmath, amssymb, amsthm, mathrsfs, color, bbm, ulem, graphicx, xy, enumerate, hyperref; the wrapper is the lane's pinned \texttt{amsart} editorial wrapper at 10pt on A4, which restates the theorem-like environments the retained text needs and adds the article's own macro definitions that the retained text needs (\texttt{\textbackslash Hom}, \texttt{\textbackslash bcdot}, \texttt{\textbackslash dD}, \texttt{\textbackslash f}, \texttt{\textbackslash gge}, \texttt{\textbackslash kk}, \texttt{\textbackslash lle}, \texttt{\textbackslash tr}), with extra wrapper packages (bm, bbm, dsfont, xspace, cancel, mathtools). The delivered numbering is the unit's own. Assets: no retained span contains a figure environment, a graphics-inclusion command, a PDF-page inclusion, a table or a TikZ picture; no image file is copied into this package, the package declares no asset dependency and no image file is redistributed with it. Licence: version arXiv:2601.22004v1 of this article is distributed under the Creative Commons Attribution 4.0 International licence, as recorded in the arXiv licence metadata for this exact version and shown on the abstract page https://arxiv.org/abs/2601.22004v1. A full rights notice is delivered at licenses/arxiv-2601.22004-CC-BY-4.0.txt. Characters: every retained excerpt is pure ASCII, so the delivered engine, XeLaTeX with its default Latin Modern text font, reports no missing character.
\begin{LEM}\label{lem_glued_t-str}
	Let $(E_1, \ldots, E_n ) $ be a full exceptional sequence in a triangulated category $\dD$ and $( F_n, \ldots, F_1 )$ its left dual. The \tr e on $\dD$ glued along $(E_1, \ldots, E_n ) $
	%filtration \eqref{eqtn_filtration} from the standard \tr es on $\dD_i/\dD_{i-1}\cong \dD^b(\kk)$ 
	is determined by the equalities:
	\begin{align*}
	\dD^{\lle 0} = \{D\in \dD\,|\, \Hom_{\dD}(D, F_s[l])=0, \textrm{ for all } l<0
	\},\\
	\dD^{\gge 0} = \{D\in \dD\,|\, \Hom_{\dD}(E_s, D[l])=0, \textrm{ for all } l<0
	\}.
	\end{align*}
\end{LEM}

\begin{align}\label{eqtn_relations}
&\beta \epsilon = 0,& &\alpha\f =0.&
\end{align}

	\begin{equation}\label{eqtn_L_S_3_P_2}
	S_3 \oplus S_3[-2] \xrightarrow{\left(\begin{array}{cc}\epsilon & \f 
		\end{array}\right)} P_2 \rightarrow L_{S_3}P_2\rightarrow S_3[1]\oplus S_3[-1].
	\end{equation}

\begin{PROP}
	The standard \tr e on $\dD^b(A)$ is glued along filtration 
	$$
	\langle S_3 \rangle  \subset \langle S_3, P_2 \rangle \subset \langle S_3, P_2, P_1 \rangle
	$$
	from the standard \tr es on $\dD^b(k)$.
\end{PROP}

\begin{proof}
	We check that the heart $A\textrm{-mod}$ of the standard \tr e $(\dD^b(A)^{\leq 0}, \dD^b(A)^{\geq 1})$ is contained in the heart of the \tr e $(\dD^b(A)^{\leq_g 0}, \dD^b(A)^{\geq_g 1})$  glued along the filtration
	$$
	\langle S_3\rangle \subset \langle S_3, P_2 \rangle \subset \dD^b(A).
	$$
	It then follows that $\dD^b(A)^{\leq 0} = \dD^b(A)^{\leq_g0}$, i.e. the \tr es coincide.  Indeed, if $A\textrm{-mod}\subset \dD^b(A)^{\leq_g0}$, then also any object of $\dD^b(A)^{\leq 0}$, admitting a `filtration' with objects of $A\textrm{-mod}[l]$ for $l\geq 0$, belongs to $\dD^b(A)^{\leq_g 0}$. Analogously, $A\textrm{-mod}\subset \dD^b(A)^{\geq_g 0}$ implies that $\dD^b(A)^{\geq 1} \subset \dD^b(A)^{\geq_g 1}$. Hence, we have an inverse inclusion of the right orthogonal subcategories $\dD^b(A)^{\leq_g 0} \subset \dD^b(A)^{\leq 0}$. 
	
	We use the description of the aisles of the glued \tr e given by Lemma \ref{lem_glued_t-str} and check that the simple $A$-modules $S_1$, $S_2$, $S_3$ lie in their intersection. Then we conlcude that $A\textrm{-mod}$, i.e. the smallest subcategory of $\dD^b(A)$ containing $S_1$, $S_2$ and $S_3$, and closed under extensions, also belongs to $\dD^b(A)^{\leq_g0} \cap \dD^b(A)^{\geq_g 0}$ as needed.
	
	As $S_3$, $P_2$ and $P_1$ are $A$-modules, clearly $S_i \in \dD^b(A)^{\geq_g 0}$, for any $i$. 
	
The sequence left dual to $\sigma$ reads
\begin{equation} \label{eqtn_left_sigma}
\langle M, L_{S_3}P_2, S_3 \rangle,
\end{equation}
where $L_{S_3}P_2$ is defined by triangle \eqref{eqtn_L_S_3_P_2}. Object $M$ is defined in two steps; as $\Hom_A(P_2,P_1) = \kk^{2}$, we have an exact triangle
\begin{equation}\label{eqtn_L_P_2P_1}
	P_2 \oplus P_2 \xrightarrow{\left(\begin{array} {cc}\alpha & \beta \end{array}\right)} P_1 \to L_{P_2}P_1 \to P_2[1] \oplus P_2[1].
\end{equation}
Applying $\Hom_A(S_3,-)$ to it and using relations \eqref{eqtn_relations}, one gets that $\Hom_A(S_3, L_{P_2}P_1) = \kk[-1] \oplus \kk[1]$. Hence, object $M$ fits into an exact triangle
\begin{equation}\label{eqtn_M}
	S_3[-1] \oplus S_3[1] \to L_{P_2}P_1 \to M \to S_3 \oplus S_3[2].
\end{equation}
As the collection \eqref{eqtn_left_sigma} is exceptional, we have $S_3\in \dD^b(A)^{\leq_g 0}$.

To calculate morphisms from $S_1$ and $S_2$ to the objects of \eqref{eqtn_left_sigma}, we use their projective resolutions
\begin{align*}
	0 \to P_2 \to P_1 \to P_3\to P_2 \oplus P_2 \to P_1 \to S_1 \to 0,&\\
	0 \to P_2 \to P_1 \to P_3 \to P_2 \to S_2 \to 0.&
\end{align*}
 The resolution of $S_1$ yields $\Hom_A^{\bcdot}(S_1, S_3 )=\kk[-2]$ and $\Hom_A^{\bcdot}(S_1, P_2) \in \dD^b(\kk)^{\geq 1}$. Hence, triangle \eqref{eqtn_L_S_3_P_2} implies that $\Hom_A^{\bcdot}(S_1, L_{S_3}P_2) \in \dD^b(\kk)^{\geq 1}$. Further, $\Hom_{A}^{\bcdot}(S_1, P_1)\in \dD^b(\kk)^{\geq 1}$, hence $\Hom_A^{\bcdot}(S_1, L_{P_2}P_1) \in \dD^b(\kk)^{\geq 0}$. Applying $\Hom_A(S_1,-)$ to \eqref{eqtn_M} one gets that $\Hom_A^{\bcdot}(S_1,M) \in \dD^b(\kk)^{\geq 0}$. It follows that $S_1\in \dD^b(A)^{\leq_g 0}$.
 
 Analogously, the projective resolution of $S_2$ implies that $\Hom_A^{\bcdot}(S_2,S_3) = \kk[-1]$ and $\Hom_A^{\bcdot}(S_2,P_2) \in \dD^b(\kk)^{\geq 3}$. It follows that $\Hom_A^{\bcdot}(S_2, L_{S_3}P_2) \in \dD^b(\kk)^{\geq 0}$.
 
 To show that negative Ext-groups from $S_2$ to $M$ vanish we first note that triangle \eqref{eqtn_L_P_2P_1} implies that object $L_{P_2}P_1$ has two non-zero cohomology objects and fits into an exact triangle
 \begin{equation}\label{eqtn_coh_L_P_2P_1}
 	S_3[1] \xrightarrow{\iota} L_{P_2}P_1 \to S_1 \to S_3[2].
 \end{equation}
This triangle implies in particular, that $\Hom_A(S_3[1], L_{P_2}P_1)$ is spanned by the truncation $\iota \colon \tau^{\leq -1}L_{P_2}P_1\to L_{P_2}P_1$. By applying $\Hom_A(S_2, -)$ to \eqref{eqtn_coh_L_P_2P_1} we see that the morphism $\Hom_A(S_2, S_3[1]) \to \Hom_A(S_2, L_{P_2}P_1)$, induced by $\iota$, is an isomorphism. As $\iota$ is the component $S_3[1] \to L_{P_2}P_1$ in \eqref{eqtn_M}, the map $\Hom_A(S_2, S_3[-1] \oplus S_3[1]) \to \Hom_A(S_2, L_{P_2}P_1)$ in the long exact sequence obtained by applying $\Hom_A(S_2, -)$ to \eqref{eqtn_M}, is an isomorphism. It follows that $\Hom_A^{\bcdot}(S_2, M) \in \dD^b(\kk) ^{\geq 0}$, hence $S_2 \in \dD^b(A)^{\leq_g 0}$. 
\end{proof}

\end{document}
